% ECONOMICS_PROBLEM: {"id": "L51-345", "title": "Privacy regulation welfare", "jel_code": "L51", "source_id": "OP-0408"}
% !TeX program = xelatex
\documentclass[11pt,a4paper]{article}
\usepackage[margin=23mm]{geometry}
\usepackage{fontspec}
\setmainfont{Latin Modern Roman}
\usepackage{amsmath,amssymb,mathtools}
\usepackage{xcolor,enumitem,booktabs,longtable,array,xurl,hyperref}
\definecolor{muted}{HTML}{53636C}
\hypersetup{colorlinks=true,urlcolor=blue,linkcolor=blue}
\setlength{\parindent}{0pt}
\setlength{\parskip}{5pt}
\newcommand{\R}{\mathbb R}
\newcommand{\E}{\mathbb E}
\newcommand{\N}{\mathbb N}
\newcommand{\ind}{\mathbf 1}
\newcommand{\Var}{\operatorname{Var}}
\newcommand{\Cov}{\operatorname{Cov}}
\newcommand{\supp}{\operatorname{supp}}
\DeclareMathOperator*{\argmin}{arg\,min}
\DeclareMathOperator*{\argmax}{arg\,max}
\newcommand{\entrymeta}[3]{{\sffamily\small\color{muted}\textbf{\texttt{#1}}\quad|\quad #2\par #3\par}}
\newcommand{\Problemref}[1]{\texttt{#1}}
\newcommand{\JELref}[2]{\texttt{#2} (JEL \texttt{#1})}
\begin{document}
\section*{Privacy regulation welfare}
\textbf{Problem ID:} \texttt{L51-345}\quad \textbf{JEL:} \texttt{L51}\par
% BEGIN ECONOMICS STATEMENT
\label{problem:OP-0408}
{\sffamily\small\color{muted}Original subject group: Digital, platform and data economics\par}
\label{definitions:problem:30007690}
\textbf{Audit classification:} Published research agenda. The nonprimary gwern mirror was replaced by an official source. Broad future research is confirmed in its indexed abstract; the original precise Section 4.6 passage and page claims were not retrievable.
\subsection*{Definitions and precise research target}
Fix European Union consumers and firms providing online retail services during a five-year period after a specified privacy-enforcement change. Define two regulatory regimes \(r\in\{0,1\}\), differing only in that enforcement policy, with their cross-border scope specified. Under either regime firms can alter data use, enter, exit and innovate. Let \(B_i(r)\) be consumer \(i\)'s monetized privacy and service benefit net of payments; let \(\Pi(r)\) be aggregate firm profit net of production, compliance and innovation costs; and let \(K(r)\) be real public enforcement cost. Define \(W(r)=\sum_i B_i(r)+\Pi(r)-K(r)\), summing over the fixed eligible-consumer population and including consumers of exiting firms. The target is \(\Delta W=W(1)-W(0)\), decomposed into privacy, service, market-structure and real compliance-cost components using a stated counterfactual ordering. Identify privacy benefits through audited data flows and informed valuation, while accounting for consent-related missing observations. Establish credible control jurisdictions or partial-identification restrictions when firms change behavior globally. Report the uncertainty associated with compliance, enforcement intensity and unobserved consumer harms. This is a dynamic welfare agenda for a specific policy change; it neither treats GDPR adoption as a clean universal experiment nor counts reduced measured tracking as a complete welfare assessment.

\textbf{Definitions, assumptions and mathematical qualifications.} Fix a monetary numeraire, eligible consumers, included firms and public agencies, and a welfare boundary. \(B_i\) includes privacy loss once; \(\Pi\) includes compliance and innovation resource costs once; fines and payments cancel only if both payer and recipient lie within that boundary. Foreign producer profits need explicit welfare weights rather than automatic cancellation. The enforcement treatment is a concrete package of audit probability, penalties, scope and implementation date, not GDPR in the abstract. To decompose \(\Delta W\), specify structural factors and a counterfactual path; correlated tracking reductions, market structure and compliance are not additively separable mechanisms by default.
\subsection*{Variants and assumption changes}
\begin{enumerate}
\item \textbf{Editorial, short-run enforcement.} Hold firm entry and innovation fixed and estimate effects of a specific enforcement change on data flows, informed user valuation and real compliance costs.
\item \textbf{Editorial, dynamic welfare boundary.} Permit global firm responses and entry, include foreign spillovers with declared weights, and compare feasible counterfactual enforcement paths under explicit equilibrium and transport assumptions.
\end{enumerate}
\subsection*{References and source audit}
\textbf{Checked source.} The Economics of Privacy, Chapter 4 (2024), pp. 97--126, primary abstract. The abstract announces a privacy-regulation research agenda. Author bibliography and indexed primary abstract checked; full official chapter returned HTTP 403. Original Section 4.6 page claims were not independently confirmed. \url{https://www.nber.org/books-and-chapters/economics-privacy/economic-research-privacy-regulation-lessons-gdpr-and-beyond}.
\begin{enumerate}\raggedright
\item\hypertarget{definitions:source:30007690:1}{} Garrett A. Johnson. 2024. \emph{Economic Research on Privacy Regulation: Lessons from the GDPR and Beyond}. The Economics of Privacy, Chapter 4 (2024), pp. 97--126, primary abstract.
\url{https://www.nber.org/books-and-chapters/economics-privacy/economic-research-privacy-regulation-lessons-gdpr-and-beyond}.
\end{enumerate}

\subsection*{Common conventions: Source mathematical conventions}

These conventions apply to each entry unless explicitly replaced.

\textbf{Models and identification.} A primitive model \(m\) specifies all probability laws, feasible choices, information, equilibrium and measurement rules used by its target. The admissible class \(\mathcal M\) is fixed independently of outcomes. If \(P_O(m)\) is its observable probability law and \(\tau(m)\) the target, its identified set at observable law \(P\) is
\[
 \mathcal I_\tau(P)=\{\tau(m):m\in\mathcal M,\ P_O(m)=P\}.
\]
Point identification means that this set is a singleton. An empty set signals incompatibility of data and assumptions. Coordinate bounds need not describe the joint set. Bounds are sharp only if every retained target is attainable or arbitrarily approachable by compatible models. Finite bounds require explicit restrictions on unobserved quantities; unrestricted confounding, hidden wealth, extrapolation or arbitrary equilibrium selection can yield uninformative sets.

\textbf{Causal interventions.} A potential outcome \(Y_i(p)\) is the outcome under a complete policy regime \(p\), including timing, financing, information and market spillovers. The target population and units are fixed before treatment. With interference, use the complete assignment vector or a justified exposure mapping; individual no-interference notation alone cannot identify an economy-wide intervention. A difference of mean potential outcomes does not require the cross-world joint distribution. Conditioning on post-treatment migration, survival, enrollment or employment changes the estimand and needs separate justification.

\textbf{Statistical statements.} An estimator, confidence set, forecast or test specifies a sampling law, model class and loss/coverage criterion. Uniform \((1-\alpha)\)-coverage means \(\inf_{m\in\mathcal M}\Pr_m\{\tau(m)\in C_n\}\geq1-\alpha\), or an explicitly asymptotic version. The whole parameter space trivially has coverage, so informativeness requires a stated width, risk or power objective. A forecasting improvement is relative to a named benchmark, information set and evaluation population.

\textbf{Equilibrium and optimization.} An equilibrium comprises feasible plans, optimality, beliefs where needed, budget constraints and all market/resource clearing. The primitives or a stated benchmark define each correspondence \(\mathcal E(p;m)\). Nonemptiness is assumed or proved, never inferred just from compact policy sets. A selected-equilibrium target uses a declared selection rule; a robust target quantifies over all permitted equilibria. Use suprema or infima unless attainment is established. An \(\varepsilon\)-optimal policy is feasible and within \(\varepsilon\) of the finite optimum value. Report an empty feasible set separately.

\textbf{Welfare and accounting.} Normative utility, interpersonal normalization, social weights, numeraire, discounting and the use of public receipts are primitives. Behavioral utility need not coincide with the welfare criterion. Household welfare includes ownership income and rents once; adding profits again to compensating variation generally double-counts them. A monetary equivalent is defined by an explicit transfer and price convention, with monotonicity and range conditions. Average effects, mechanism contributions, transfers and welfare losses are different targets. Infinite sums require convergence and dynamic feasible plans require terminal or no-Ponzi conditions.

\textbf{Source versions.} A working-paper date, revision date and journal publication year are distinguished. A DOI for a journal version is not automatically the DOI of an earlier manuscript. Locators refer to the stated version and printed pages where specified. A metadata-only check does not verify a claim about a full-text conclusion. Every entry's source subsection narrows the provenance claim accordingly.
% END ECONOMICS STATEMENT
\end{document}
